\section{How few links are needed}\label{sec:minimal}

How small can a network be and still hide a regime change in this way? The answer depends on the cost class and on the type of breakpoint. Table~\ref{tab:minimaledges} collects what is proved. Throughout, demands are proportional with every $r_k>0$, every simple path is a route, and ``equal slopes'' means $E'(t_0^-)=E'(t_0^+)$ at a nested breakpoint $t_0>0$.

\begin{theorem}[Minimal link counts]\label{thm:minimal}
Consider networks with costs $a_e+b_ez^p$, $a_e,b_e>0$, and one common real exponent $p>1$. Every network with at most four edges has strict slope jumps of $E$ and of $P$ at every nested breakpoint, in the direction of Theorem~\ref{thm:weakcomparison}. Hence five is the minimum number of edges permitting equal slopes, both at nested expansions and at nested contractions.
\end{theorem}

The lower bound is Theorem~\ref{thm:fouredge}, proved in Appendix~\ref{app:minimality}. The five-link expansion and contraction of Section~\ref{sec:fiveedge} attain it. The proof is a finite case analysis with an explicit inequality closing each case (Table~\ref{tab:minimality}). Its outline is as follows.
\begin{itemize}[nosep]
\item By Lemma~\ref{lem:tangencycomparison}, equal slopes force a tangency of the gap of a new route.
\item If the smaller regime has no demand-preserving load changes ($\Kc_{R_s}=\{0\}$), the common-power identity $Dq=p(c_0-a)/t_0$ makes the gap a function $\alpha+\beta t^p$, which cannot be tangent to zero without vanishing identically. This disposes of every connected graph with at least four vertices or at most three edges.
\item Networks whose cycles are all parallel banks have no contractions, and their expansions are strict because a joining branch lowers the derivative of the bank delay to zero.
\item The remaining graph is a transitive triangle with one side doubled. There, a curvature inequality for parallel banks, $\ell''-(p-1)\ell'/t>0$ unless all free-flow costs in the bank are equal, rules out every tangency.
\end{itemize}

\begin{table}[htbp]
\centering\small
\begin{tabularx}{\linewidth}{@{}>{\raggedright\arraybackslash}p{4.3cm}>{\raggedright\arraybackslash}p{2.3cm}>{\raggedright\arraybackslash}p{2.6cm}X@{}}
\toprule
Cost class & Breakpoint & Minimum edges & Evidence\\
\midrule
Strictly convex polynomials & expansion & $3$ & Example~\ref{ex:cubic}; two edges never suffice (Remark~\ref{rem:twoedges})\\
BPR, edge-specific exponents & contraction & $3$ & Example~\ref{ex:hetero}; Remark~\ref{rem:twoedges}\\
BPR, edge-specific exponents & expansion & $4$ or $5$ (open) & three edges excluded (Proposition~\ref{prop:hetero}, Remark~\ref{rem:hetero}); five-link example\\
BPR, common exponent & expansion & $5$ & Theorems~\ref{thm:main} and~\ref{thm:fouredge}\\
BPR, common exponent & contraction & $5$ & bank-swapped example; Theorem~\ref{thm:fouredge}\\
\bottomrule
\end{tabularx}
\caption{Minimum numbers of edges for equal one-sided slopes at nested breakpoints. Positive free-flow and congestion coefficients throughout.}
\label{tab:minimaledges}
\end{table}

\begin{remark}[Two edges never suffice]\label{rem:twoedges}
Let every cost be $C^1$ with $c_e'(z)>0$ for $z>0$. If a network with at most two edges has a breakpoint, some commodity has two routes, so the two edges are parallel, say $s\to v$, and carry a total demand $Qt$ with $Q>0$. Their common delay $\ell(Qt)$ is strictly increasing in $t$, so an edge can only join: there are no contractions. Suppose edge $2$ joins at $t_0$, so $c_2(0)=\ell(Qt_0)$. To the left only edge $1$ is used, with positive load, and $\ell'(Qt_0^-)=c_1'(Qt_0)>0$. To the right the loads satisfy $c_1(y_1)=c_2(y_2)$ with $y_1+y_2=Qt$. Hence $y_2/(Qt-Qt_0)\to c_1'/(c_1'+c_2'(0))$, which is read as $1$ if $c_2'(0)=0$, and $\ell'(Qt_0^+)=c_1'c_2'(0)/(c_1'+c_2'(0))<c_1'$. Since $E=Qt\,\ell(Qt)$, the slope of $E$ drops strictly, by $Q^2t_0[\ell'(Qt_0^-)-\ell'(Qt_0^+)]>0$.
\end{remark}

These minima concern the stated classes only. In the two-bank architecture a sharper, order-dependent count holds: contacts of orders $m_1,\dots,m_\ell$ need exactly $2+\sum_jm_j$ edges (Appendix~\ref{app:contacts}). That count is specific to the architecture and is not a lower bound for arbitrary graphs.
